
\begin{figure}[t]
\centering
\includegraphics[width=0.92\linewidth]{figures/residual_chain.pdf}
\caption{The structural-residual reduction chain (\vref{v175}~$\to$ the certification round). The naive
``build a quantum-gravity measure'' shrinks, one machine-checked step at a time, from a vague measure to
one falsifiable physical statement --- \emph{is the raw seam $(E_8)_1$ at $\tau=i$?}
(\texttt{SEAM.EQUIV.01}; its conformal-deck face \texttt{QGEO.SYM.01} is a corollary, \vref{v335}).
The closing arc (\vref{v276}--\vref{v302}) pins the shared fact and derives the gap $6\ln(3/2)>0$; the
certification round discharges the extension leg to the crossed-product package (\vref{v469}), the R3
graph$\to$geometry bridge to Kronheimer 1989 (\vref{v479}) and exhibits the four-interval $\mu_4$
mechanism (\vref{v480}). Every step above the bedrock is a theorem, a Lean proof or an established
citation; the residual --- the cited continuum scaling-limit existence (\vref{v336}) --- stays honestly \mO.}
\label{fig:residual-chain}
\end{figure}

\begin{keybox}[\texttt{QGEO.SYM.01} is \emph{not} a theorem --- it is the final seam-identification premise]
\textbf{The single load-bearing bolt of the whole construction, on one page.}
\begin{enumerate}[leftmargin=1.6em,itemsep=1pt]
\item \textbf{What is identified.} That the conformal structure of the raw RP seam double \emph{is} the
$\mu_4$-deck structure of the carrier clock --- i.e.\ the carrier (P2) and the seam's conformal geometry
are the same object. Equivalently, in Riemann--Roch \emph{index} language: the seam's degree-$4$ mark
divisor on the genus-$0$ double has $h^0=\deg{+}1=5=\gcar$ and $\operatorname{rank}H_1=\deg{-}1=3=\Nfam$
(with $\Lambda^{\mathrm{even}}(\mathbb C^5)=1{+}10{+}5=16=\dim S^+$), so P2 is the genus-$0$ mode count of the
four-mark seam divisor (\veri{v228\_rr\_index\_gate.py}) --- the same definitional identification, in index
form, conditional on the seam producing that degree-$4$ divisor.
\item \textbf{Why it is weaker than \texttt{QGEO.CONF.01}/\texttt{ISO.01}.} It does \emph{not} name
$\PP$ or $\mu_4$ (uniformisation \emph{produces} them) and asks only for a conformal \emph{symmetry},
not a metric isometry or a pre-named conformal class --- strictly the mildest of the three.
\item \textbf{What follows automatically (the conditional theorem).} \emph{Given} \texttt{QGEO.SYM.01},
equivariant uniformisation puts the order-$4$ clock in the form $z\mapsto iz$ on $\PP$, the four
parabolic marks are its free orbit $\mu_4$, $H^1$ carries the $(1,2,3)$ grading, the seam Calder\'on
contraction is the rank-$8$ $K_\Sigma$ polarisation, and the index-$4$ simple-current extension is
$(E_8)_1$ --- the \emph{entire} AQFT closure of \vref{v175}--\vref{v181}.
\item \textbf{What would falsify or unseat it.} A seam double of genus $>0$ (uniformisation to $\PP$
fails); a carrier clock that is \emph{not} a conformal symmetry of the seam geometry (a non-deck
transport); or a fourth chiral generation ($\Nfam\neq3$ collapses $b_1=\Nfam$ and the four-mark $\mu_4$
structure). The first two are constructive-geometry questions about the raw seam; the third is the
standard family-count kill test.
\end{enumerate}
\noindent\textbf{Honest framing.} The correct external statement is \emph{conditional}: ``\emph{Given}
\texttt{QGEO.SYM.01}, the RP seam boundary realises the $\mu_4$ deck structure, and the AQFT closure to
$(E_8)_1$ follows through the verified chain \vref{v175}--\vref{v181}'' --- \emph{not} ``the theory
\emph{derives} the seam geometry''. The premise stays \mO; everything downstream of it is \mE. Per
\vref{v323} (Bisognano--Wichmann) \texttt{QGEO.SYM.01} is itself the \emph{downstream} conformal-deck face
of the single named keystone \texttt{SEAM.EQUIV.01} (the raw RP seam $=$ the holomorphic $(E_8)_1$ net at
$\tau{=}i$), which is now \emph{closed modulo cited theorems on its MMST route} \texttt{SEAM.EQUIV.MMST.01}
(the lattice/$S_3$ stack pins it at every computable
level, Lean \texttt{FORM.SEAM.MMST.01}; only the cited continuum existence stays \mO; the twistor route
\texttt{SEAM.EQUIV.TWISTOR.01} and the unconditional parent stay \mO).
\par\smallskip\noindent\textbf{This is not a gap to apologise for --- it is TFPT's one fundamental
postulate.} Every fundamental theory rests on an irreducible premise (the constancy of $c$ in
relativity, the equivalence principle, the canonical commutator). \texttt{QGEO.SYM.01} is that premise
for TFPT: a single geometric identification (seam conformal structure $=$ carrier $\mu_4$ deck) from
which the entire algebraic and topological structure unfolds \emph{by theorem}. The achievement is not
that there is no postulate --- it is that the whole theory has been compressed onto \emph{exactly one},
sharply located and explicitly falsifiable.
\end{keybox}

\begin{keybox}[The Seam-Holomorphy Theorem --- the one named closing target \mE/\mO\ (\veri{v234\_seam\_holomorphy\_selection.py})]
The bedrock above can be stated as \emph{one condition with three provably-equivalent faces}, so that
``closing the theory'' becomes a single, sharply-bounded analytic theorem. \textbf{The condition:} the seam
carries \emph{no nontrivial abelian sector}. Its three faces all force $E_8$:
\begin{enumerate}[leftmargin=1.7em,itemsep=1pt]
\item \textbf{AQFT} (\vref{v83}/\vref{v154}): the boundary chiral net is \emph{holomorphic} ($\mu$-index $1$,
single vacuum sector). Holomorphic $+\,c=8$ $\Rightarrow$ the unique even unimodular rank-$8$ lattice net
$(E_8)_1$, and $c=8=\gcar+\Nfam$ is itself structural.
\item \textbf{Geometry} (\vref{v232}): the seam link is a \emph{homology $3$-sphere} $\Leftrightarrow$ its
binary group $\Gamma$ is \emph{perfect} $\Leftrightarrow$ $\Gamma=2I$ (the Poincar\'e sphere $S^3/2I$).
\item \textbf{Representation theory} (\vref{v219}): exactly \emph{one} $1$-dim irrep (the trivial one).
\end{enumerate}
These are the \emph{same} number: $\#(1\text{-dim irreps})=|\Gamma^{\mathrm{ab}}|=|H_1(S^3/\Gamma)|
=\#(\text{mark-}1\text{ affine-Dynkin nodes})$, which equals $1$ \emph{only} for $E_8$ ($A_n{\to}n{+}1$,
$D_n{\to}4$, $E_6{\to}3$, $E_7{\to}2$, $E_8{\to}1$). So $2I$ is the unique nontrivial perfect finite
$SU(2)$ subgroup and $S^3/2I$ the unique homology-sphere space form. \textbf{Consequence:} Research
Contract~2 ($G_{\mathrm{net}}$/Target~A), the carrier P2, and the Kleinian seam are \emph{not} separate
gates --- they are this \emph{one} condition.
\par\smallskip\noindent\textbf{The closing chain (what a complete closure would be), \mO\ at one step:}
\[
\begin{gathered}
  \text{RP}+\text{gap}+\text{chirality}\ \xrightarrow{\ \vref{v160}\ }\ \text{quasi-free bulk}
  \ \xrightarrow{\ ?\ }\ \text{holomorphic boundary net}\\[2pt]
  \xrightarrow{\ \mE\ }\ (E_8)_1\ \xrightarrow{\ \vref{v219}/\vref{v228}\ }\ 2I,\ (2,3,5),\ \text{P2}
  \ \xrightarrow{\ \vref{v179}\text{--}\vref{v181}+\text{Kleinian genus-}0\ }\ \texttt{QGEO.SYM.01}.
\end{gathered}
\]
The free-bulk premise is already a fixed-point theorem (\vref{v160}--\vref{v165}); the Kleinian model
supplies the genus-$0$ of the uniformisation chain for free. \textbf{The single remaining analytic step is
``free Gaussian RP bulk $\Rightarrow$ holomorphic (single-sector) boundary net''} --- the bulk--boundary
statement that an invertible (Gaussian) bulk has a holomorphic chiral boundary. That step is now \emph{closed
modulo cited theorems} (the MMST route \texttt{SEAM.EQUIV.MMST.01}): the lattice/$S_3$ stack pins it at every
computable level (\vref{v367}/\vref{v368},
\vref{v376}--\vref{v379}, Lean \texttt{FORM.SEAM.MMST.01}), only the cited continuum existence staying \mO\
(\vref{v336}); $v_{\mathrm{geo}}$ stays the No-Unit metrology primitive and
$F_{\mathrm{transfer}}$ stays external physics. This box certifies that the three faces coincide and
force $E_8$ (\mE); the analytic step is closed modulo that cited theorem.
\par\smallskip\noindent\textbf{A physical realisation of the step (\veri{v235\_seam\_chern\_simons.py}).}
The step lands in standard physics. A free gapped \emph{bosonic} $2{+}1$d bulk is an abelian
Chern--Simons / $K$-matrix theory (even integer $K$), where $\#$anyons $=|\det K|$, the edge chiral central
charge is $c={\rm signature}(K)$, and the edge net is \emph{holomorphic} $\Leftrightarrow|\det K|=1$ --- so
``no nontrivial abelian sector'' is precisely $\det K=1$. The extension tower of \vref{v92} \emph{is} the
anyon-condensation tower read by determinant: $D_5{\oplus}A_3$ ($\det16$, $16$ anyons) $\to SO(16)_1{=}D_8$
($\det4$) $\to (E_8)_1$ ($\det1$, holomorphic), all rank $8$ with $c={\rm signature}=8=\gcar+\Nfam$, so the
holomorphic end is the \emph{Kitaev $E_8$ quantum-Hall state}. Holomorphy $\Leftrightarrow$ condensing the
order-$|\mu_4|=4$ \emph{Lagrangian} glue ($\det16{\to}16/4^2{=}1$); a partial order-$2$ isotropic only
reaches $\det4$ ($D_8$). The residual is thereby a single integer condition --- ``the free RP seam condenses
the order-$|\mu_4|$ Lagrangian glue'' --- the same sheet/$\mu_4$ selection as \texttt{GATE.QGEO}, but now with
$\text{holomorphic}\Leftrightarrow\det1$ a theorem (\texttt{GATE.HOLO.02}, \mE/\mO).
\par\smallskip\noindent\textbf{The residual as a \emph{physical} statement (\veri{v237\_seam\_sre\_closure.py}).}
The same integer condition is the physics of topological order: a $K$-matrix phase has ground-state
degeneracy $|\det K|^g$ on a genus-$g$ surface, so $\det K=1$ $\Leftrightarrow$ \emph{no} topological
ground-state degeneracy on any closed surface $\Leftrightarrow$ the bulk is \emph{short-range-entangled}
(invertible). Among rank-$8$ even positive-definite $K$ ($c=8=\gcar+\Nfam$) only $E_8$ is SRE --- the unique
nontrivial bosonic SRE $2{+}1$d phase, the Kitaev $E_8$ state, whose edge is the holomorphic $(E_8)_1$. TFPT
already has a unique vacuum on the plane (RP $+$ gap $+$ cluster $+$ OS); the plane cannot see $\det K$ (no
non-contractible cycles), the torus can. So the one analytic step is the sharp, falsifiable physical
statement \emph{``the free RP gapped seam bulk is short-range-entangled (no topological order)''}
($=\det K=1=$ the Kitaev $E_8$ phase, \texttt{GATE.HOLO.03}, \mE/\mC) --- and the lattice/$S_3$ stack now answers
it at every computable level (torus $\det K{=}1$, \vref{v378}), so the keystone's MMST route
(\texttt{SEAM.EQUIV.MMST.01}) is \emph{closed modulo cited
theorems}, only the cited continuum existence (\vref{v336}) staying \mO\ (the parent stays \mO).
\end{keybox}

\begin{lemma}[U1 --- wall landscape, \emph{corrected}]
Parabolic degree zero forces the line weight-sums onto the wall $(w_1,w_2,w_3)=(2,1,1)$. An explicit
balanced representative (weights in $\tfrac13$) is
\[
  W_{\mathrm{wall}}=\tfrac13\begin{pmatrix}2&1&2&1\\1&0&0&2\\0&2&1&0\end{pmatrix},
\]
each column a permutation of $(0,1,2)$, row sums $(2,1,1)$. \textbf{Honest result:}
the $6^4{=}1296$ column-permutation matrices give $144$ walls; under the \emph{full} group
$D_4\times S_3(\text{lines})\times S_3(\text{values})$ they form \emph{five} orbits, \emph{not one}. So
$W_{\mathrm{wall}}$ is \emph{not} singled out by symmetry --- the uniqueness must come from the selector
(Lemma U4), not from a symmetry quotient. The wall landscape is nonetheless a finite, explicitly
enumerated set; $W_{\mathrm{wall}}$ is one of the five. Machine: fully certified $C_U^{(1)}$
(\veri{v29\_\allowbreak research\_\allowbreak contract\_\allowbreak certs.py}).

\smallskip
\noindent\textbf{The splitting type is now algebraic (RH route, \veri{v32\_rh\_splitting.py}).} Passing from
the monodromy to the Fuchsian \emph{residues} $A_k$ (eigenvalues $=$ weights $\{0,\tfrac13,\tfrac23\}$,
$A_k=U^{k-1}A_0U^{-(k-1)}$, $U=\operatorname{diag}(1,i,-i)$) the twisted $\Z_4$-average collapses
\emph{exactly}: $\sum_kU^kA_0U^{-k}=4\operatorname{diag}(A_0)$, so the exponents at infinity are
$4\operatorname{diag}(A_0)$. A flat bundle with regular $\infty$ needs integer exponents summing to
$4{=}-\deg E$; with the cusp weights the only options are perms of $(2,1,1)$ and $(2,2,0)$, i.e.
\[
  \boxed{\ \mathcal O(-2)\!\oplus\!\mathcal O(-1)^2\ \Longleftrightarrow\ \operatorname{diag}(A_0)=(\tfrac12,\tfrac14,\tfrac14)\ }
\]
(the sibling $(2,2,0)\to\mathcal O(0)\oplus\mathcal O(-2)^2$). Such a Hermitian $A_0$ (spectrum
$\{0,\tfrac13,\tfrac23\}$, that diagonal) exists by Schur--Horn. So the \emph{splitting} is an algebraic
condition on $\operatorname{diag}(A_0)$. \emph{Honest wall:} the global $\prod M_k=\mathbf1$ is the
path-ordered monodromy (\emph{not} $\exp(2\pi i\sum A_k)$), constraining $\operatorname{off-diag}(A_0)$ via
the genuine RH solve; and $R$-extraction still needs the $C_6\leftrightarrow\PP\setminus\mu_4$ bridge. The
RH route fixes the \emph{splitting} algebraically but not $c_u/c_d$.
\end{lemma}

\begin{lemma}[U2 --- the kill switch, \emph{run and hardened} \mE]
The selected polystable wall representative must \emph{not} collapse to the diagonal abelian $U(1)^3$
point if it is to give non-trivial $SU(3)_F$ transport mixing: \textbf{(A)} $\nabla_F^\star$ polystable but
\emph{not} simultaneously diagonal (the desired breakthrough), or \textbf{(B)} diagonal collapse, in which
case $c_u/c_d$ stays honestly \mC. \textbf{Result: the cheap test is run robustly and lands on A.} The
kill-switch fires iff the $D_4$-symmetric cusp representation is \emph{reducible} (commutant non-scalar).
Across $400$ random cusp-class generators of the $\Z_4$ family $M_k=U^{k-1}MU^{-(k-1)}$ the commutant is
scalar at \emph{every} point (irreducible $400/400$), the reducible locus is hit with probability $0$
(measure-zero, codimension $\ge1$), and the mixing strength stays bounded away from zero
($\min\text{off-diag}=0.35>0.05$). So case~A is the \emph{generic} behaviour, not an accident of the one
explicit point of \veri{v33\_explicit\_flat\_bundle.py}: $(U_{\mathrm{wall}})$ \emph{survives} its own
cheapest falsification test. \veri{v38\_uwall\_killswitch.py}
\end{lemma}

\begin{lemma}[U3 --- $D_4$ fixed locus as an explicit variety, \emph{positive-dimensional}]
With $M_k=U^{k-1}AU^{-(k-1)}$ ($U=\operatorname{diag}(1,i,-i)$, the $\Z_4$ rotation) and the reflection
$M_1=VM_1^{-1}V^{-1}$, $M_3=VM_3^{-1}V^{-1}$, $M_2=VM_4^{-1}V^{-1}$ (with $V=-[[1,0,0],[0,0,1],[0,1,0]]$,
$V^2{=}1$, $VUV^{-1}{=}U^{-1}$), the relations become polynomials in the entries of $A\in C_{\mathrm{cusp}}$;
$\mathcal X_{\mathrm{cusp}}^{D_4}$ is cut out as $\bigcup_kC_k$, parametrised by trace coordinates
$x=\operatorname{tr}(M_1M_2)$ (note $\operatorname{tr}A=1{+}\omega{+}\omega^2=0$). \textbf{Result
($C_U^{(2)}$, \veri{v30\_d4\_character\_variety.py}):} adding the reflection to the $\Z_4$ locus of \vref{v19} does
\emph{not} isolate the point --- the \emph{full} $D_4$-fixed product locus is still
\emph{positive-dimensional} ($|\operatorname{tr}(M_1M_2)|$ varies continuously over $\approx[0,3]$). So even
the full $D_4$ symmetry cannot select $\nabla_F^\star$. \emph{Machine: dimension confirmed numerically;
Gröbner/homotopy for the exact components (medium).}
\end{lemma}

\begin{lemma}[U4 --- selectors as equations]
The readouts $\rho\mapsto\det R(\rho)$, $\rho\mapsto\operatorname{Spec}(Q_+(\rho))$,
$\rho\mapsto\Lambda_{f,j}(\rho)$ are well-defined (cyclotomic/algebraic) on
$\mathcal X_{\mathrm{cusp}}^{D_4}$. \textbf{Acceptance criterion:}
\[
  \#\{\rho\in\mathcal X_{\mathrm{cusp}}^{D_4}:\ \det R=8,\ \operatorname{Spec}(Q_+)=\{1,2,3\}\}=1
\]
in the polystable quotient. If more than one survives, the theory needs a further input and the gate is
\emph{not} closed. \textbf{The bottleneck (made precise by U3):} since the $D_4$-fixed variety is
positive-dimensional, the \emph{selector} must do all the cutting --- and evaluating $\det R(\rho)$,
$\operatorname{Spec}(Q_+(\rho))$ as functions on $\mathcal X_{\mathrm{cusp}}^{D_4}$ requires the
\emph{parabolic-degree\,$\leftrightarrow$\,residue dictionary} $R(\rho)$, i.e.\ the H2 equivalence. That
dictionary was the single remaining analytic input of the $(U_{\mathrm{wall}})$ gate. \emph{Status update:
since \vref{v118}--\vref{v122} the dictionary itself is exact} --- hexagon $=$ sign-twisted family spectrum, lepton
coefficients $=$ resolvent determinants of the $W(A_3)$ monodromy, addresses pinned, margins theorems
(boxes below); what remains is the strictly smaller R4$'$, the \emph{establishment of the selectors}
($n=(5,-9,6)$ $+$ the diamond determinants; since \veri{v134\_dual\_anchor.py}: the pair $(d,n)$).

\textbf{What $R(\rho)$ is (\veri{v31\_R\_dictionary.py}).} Two facts pin its nature. \emph{(i) Case A is
alive:} every $D_4$-fixed cusp monodromy found is \emph{irreducible}, so the wall point carries
non-trivial $SU(3)_F$ mixing --- the U2 kill switch does \emph{not} fire to the diagonal case B.
\emph{(ii) $R(\rho)$ is not algebraic:} $R$ is integer-valued, but the algebraic invariants of $\rho$
(e.g.\ $\operatorname{tr}(M_1M_2)$) vary \emph{continuously} on the locus; an integer readout cannot be a
continuous algebraic function. Hence $R(\rho)$ is the discrete invariant of the \emph{non-abelian Hodge /
Mehta--Seshadri} map (the parabolic-degree / Hodge-filtration jumps of the harmonic bundle $E_\rho$):
integer output, but its \emph{computation} is the transcendental harmonic-metric (Hitchin) solve \emph{on
the generic Higgs locus}. \textbf{Refinement (U4$''$, \veri{v40\_harmonic\_metric.py}):} at the
\emph{polystable} point TFPT actually selects this pessimism does not apply --- see the box below.
\end{lemma}

\par\smallskip\noindent\textbf{Progress: the \emph{selector side} of U4 is closed on the explicit point} (\mE).
While $R(\rho)$ as a function on the whole locus needs the Hitchin solve, the two U4 \emph{selectors} are
\emph{algebraically} accessible on the explicit flat bundle of \veri{v33\_explicit\_flat\_bundle.py} --- no
PDE. Computed exactly from $A_0$: the exponents at infinity $\operatorname{eig}(\sum_kA_k)=\{2,1,1\}$ give
the splitting type $E=\mathcal O(-2)\oplus\mathcal O(-1)^2$, i.e.\ the parabolic \emph{anchor} $a=(1,1,2)$;
the parabolic degree is $0$ ($\deg E=-4=-|\mu_4|$ plus four cusp weight-sums of $1$), so the point is
\emph{polystable} (Mehta--Seshadri substrate). The two selectors then read off the bundle directly:
\[
  \boxed{\ \begin{gathered}
  \det R=8=n\!\cdot\!a\quad(\text{anchor pairing on the splitting type}),\\[2pt]
  \operatorname{Spec}(Q_+)=\{1,2,3\}=3\alpha+1\quad(\alpha=\text{cusp weights }\{0,\tfrac13,\tfrac23\})
  \end{gathered}\ }
\]
So $\det R=8$ is the anchor pairing of the splitting type and $\operatorname{Spec}(Q_+)$ is the affine image
($d{=}3$ the weight denominator) of the cusp weights --- \emph{both selectors are read off the explicit
bundle's algebraic data}. \textbf{Residual (sharpened below by Readout Rigidity):} the quark \emph{ratio}
$c_u/c_d$ turns out \emph{not} to need the Hitchin datum at all --- it is \emph{constant} on this discrete
selector stratum (\vref{v49}), hence an integer Pl\"ucker readout. The only genuinely transcendental input left is
the \emph{absolute} amplitude normalisation ($U_{\mathrm{point}}$), which is an anchor, not a missing PDE.
\veri{v39\_uwall\_selectors.py}


\par\smallskip\noindent\textbf{Breakthrough: the harmonic metric is \emph{finite linear algebra}, not a PDE} (\mE).
The U4 ``transcendental Hitchin solve'' is the worst case for a \emph{generic} Higgs bundle, but the point
TFPT selects is \emph{polystable}, and by Mehta--Seshadri a parabolic-degree-$0$ polystable bundle is
\emph{unitary} --- a unitary representation has the \emph{constant} invariant Hermitian harmonic metric and
zero Higgs field $\Phi=0$. So at the selected point the harmonic metric is finite linear algebra. On the
explicit bundle of \veri{v33\_explicit\_flat\_bundle.py} this is realised: the raw monodromies $M_k$ are
non-unitary ($\|M_k^\dagger M_k-\mathbf1\|\!\approx\!0.83$), but there is a \emph{unique} common invariant
Hermitian form $H$ ($\dim\ker\{M_k^\dagger HM_k=H\}=1$), it is \emph{positive-definite}, and the unitarised
holonomy $\widetilde M_k=H^{1/2}M_kH^{-1/2}$ is unitary ($\sim10^{-8}$), lies in the cusp class
$\{1,\omega,\omega^2\}$ with $\prod\widetilde M_k=\mathbf1$, and has \emph{clean} matrix elements
$|\widetilde M_k|_{ij}\in\{0,\tfrac12,\tfrac1{\sqrt2}\}$. \textbf{So $C_U^{(3)}$ reduces from a transcendental
PDE to: this finite unitarisation (done) $+$ the combinatorial word-length $R$ (closed: H1 $+$ $\det R=8$).}
\emph{Residual:} the feared PDE is gone; the quark \emph{ratio} $c_u/c_d$ is then fixed combinatorially by
Readout Rigidity (below), and only the \emph{absolute} amplitude scale remains as an anchor.
\veri{v40\_harmonic\_metric.py}


\par\smallskip\noindent\textbf{The final leg-assignment test (honest result)} (\mE).
Does the harmonic-frame holonomy $\{0,\tfrac12,\tfrac1{\sqrt2}\}$ feed the $\delta=\tfrac12$ resolvent to
give the lepton amplitudes $\Lambda=(0.475,1.107,0.917)$? \textbf{(i) Ruled out as a literal identity:}
$\Lambda_\mu=1.107>1$, while unitary holonomy entries have modulus $\le1$, so the amplitudes \emph{cannot}
be holonomy matrix elements --- they are the non-unitary resolvent Green function (as \vref{v19} already noted).
\textbf{(ii) Suggestive positive:} the harmonic-frame holonomy diagonal modulus is exactly
$|\mathrm{diag}\,\widetilde M_0|=(0,\tfrac12,\tfrac12)$, and that $\tfrac12$ \emph{equals} the distinguished
lepton transport value $\delta=\tfrac12$ that \vref{v20} used by hand --- so the geometry plausibly \emph{supplies}
$\delta$ (an observation at the explicit point; genericity across the locus is future work).
\textbf{(iii) Clean negative:} no natural holonomy construction (e.g.\ $(\mathbf1-\tfrac12\widetilde M)^{-1}$)
maps the holonomy to the amplitude vector. \textbf{Verdict:} the amplitudes are the $\delta=\tfrac12$
resolvent (closed combinatorially via the word-length $R$); the holonomy supplies $\delta$ and the leg
selection, not the amplitude magnitude. $c_u/c_d$ is \emph{not} obtained from the holonomy alone --- a clean,
honest negative on the literal reproduction, with one suggestive geometric link ($\delta=\tfrac12$). The
negative is now \emph{explained}: the scalar leg is wrongly typed --- see the Exterior Leg Lemma.
\veri{v41\_leg\_assignment.py}


\par\smallskip\noindent\textbf{The torsion normal form and the $\delta=\tfrac12$ theorem} (\mE).
The \vref{v41} ``genericity is future work'' item is now executed --- and yields a theorem. \emph{(i) Flatness
$=$ $\mu_4$ torsion \mE:} for \emph{any} $M$ (symbolic), $\prod_k U^kMU^{-k}=(MU)^4$ --- the $\Z_4$-family
flatness constraint \emph{is} the torsion statement ``$T=MU$ is a fourth root of unity'': the $\mu_4$
atom appears as the \emph{order} of the twisted cusp generator. \emph{(ii) The $\delta=\tfrac12$ theorem
\mE\ (both directions):} on the \emph{involutive} branch ($T^2=\mathbf1$) $T$ is a Hermitian reflection
$2vv^\dagger-\mathbf1$, and the cusp trace condition $2v^\dagger U^{-1}v=1$ splits \emph{exactly} into
$|v_1|^2=\tfrac12$ (real part) and $|v_2|=|v_3|$ (imaginary part), forcing
\[
  \boxed{\ \operatorname{diag}M=(0,\tfrac i2,-\tfrac i2),\qquad \operatorname{spec}M=\{1,\omega,\omega^2\}\
  \text{automatic}\ }
\]
($\chi_M=\lambda^3-1$ from unitary $+$ trace $0$ $+$ $\det 1$). So $\delta=|M_{22}|=\tfrac12$ holds
\emph{exactly and constantly} on the whole two-parameter branch --- the distinguished lepton transport
value that \vref{v20} used by hand is a \emph{$\mu_4$-torsion identity}, not an observation. \emph{(iii)
Reflection lemma \mE:} the $D_4$ reflection forces $\operatorname{diag}M=(r,z,\bar z)$ with $r$ real and
$\operatorname{Re}z=-r/2$ --- the whole $D_4$ diagonal freedom is \emph{one} real parameter, and
$\delta=\tfrac12\iff r=0$. \emph{(iv) Branch census \mE\ (seeded):} the full $D_4$-fixed flat cusp locus
realises only $\operatorname{tr}(MU)\in\{1,-1\}$; the involutive branch has the diagonal \emph{constant}
$(0,\tfrac12,\tfrac12)$ to machine precision, while on the $\{1,i,-i\}$ branch $d_1$ varies over $[0,1]$
--- and the explicit \vref{v33}/\vref{v40} point sits exactly in its $d_1=0$ slice, the \emph{same} $\delta$ value.
\emph{Residue \mC\ (recorded):} whether the anchor splitting $\mathcal O(-2)\oplus\mathcal O(-1)^2$
(Mehta--Seshadri side) forces the $d_1=0$ slice --- or selects the involutive branch --- is the remaining
question; $\delta=\tfrac12$ itself is no longer the open item. \veri{v114\_torsion\_delta.py}


\par\smallskip\noindent\textbf{The anchor pins the residue: an exact normal form for the $(U_{\mathrm{wall}})$ bundle} (\mE).
The \vref{v114} branch question is answered --- and the \vref{v33} \emph{numerical} Riemann--Hilbert point acquires an
\emph{exact closed form}. \emph{(i) $\mu_4$-average lemma \mE:} $\sum_kU^kXU^{-k}=4\operatorname{diag}X$
for any $X$ --- $\mu_4$ conjugation-averaging \emph{is} the diagonal projection, so the exponents at
infinity are $4\operatorname{diag}A_0$ and the anchor splitting holds \emph{iff}
$\operatorname{diag}A_0=(2,1,1)/4=a/|\mu_4|$: the residue diagonal \textbf{is} the anchor over $\mu_4$
(\vref{v33}'s diagonal was forced, not chosen). \emph{(ii) The $(8,0,5)/144$ lemma \mE:} the anchor diagonal
and the cusp spectrum $\{0,1,2\}/\Nfam$ pin $\sum|a_{ij}|^2=\tfrac{13}{144}$, and with $a_{13}=0$ the
determinant condition solves \emph{uniquely} to
\[
  \boxed{\ \bigl(|a_{12}|^2,\,|a_{13}|^2,\,|a_{23}|^2\bigr)=\frac{(8,\,0,\,5)}{144}\ }
  \quad\begin{aligned}
  &\scriptstyle 8=\operatorname{rank}E_8,\ 5=\gcar,\\[-2pt]
  &\scriptstyle 144=(|\mu_4|\Nfam)^2,\ 8+5=13=\Delta_Q
  \end{aligned}
\]
--- the off-diagonal weights are compiler atoms, and the total $13$ is the \emph{same} $\Delta_Q$ that
denominates $c_u/c_d=55/117=5\cdot11/(9\cdot13)$. \emph{(iii) Exact normal form \mE:}
$A_0^\star=\bigl[\begin{smallmatrix}1/2&\sqrt2/6&0\\\sqrt2/6&1/4&\sqrt5/12\\0&\sqrt5/12&1/4\end{smallmatrix}\bigr]$
has $\chi=\lambda(\lambda-\tfrac13)(\lambda-\tfrac23)$ \emph{exactly}. \emph{(iv) $A_0^\star$ is the RH
solution \mE:} its big-circle monodromy is trivial to $10^{-10}$, $\prod M_k=\mathbf1$, and the
unitarised holonomy has $|\operatorname{diag}\widetilde M_0|=(0,\tfrac12,\tfrac12)$,
$\operatorname{tr}(\widetilde M_0U)=1$ --- the \vref{v33}/\vref{v40} point is gauge-equivalent to $A_0^\star$ (the
$a_{12},a_{23}$ phases are diagonal gauge). \emph{(v) The anchor forces the $d_1=0$ slice \mE:} a
multi-seed RH scan over the anchor-pinned spectral manifold lands at \emph{every} flat solution on the
\emph{same} gauge orbit with $|\operatorname{diag}\widetilde M_0|=(0,\tfrac12,\tfrac12)$ --- combined
with \vref{v114}, the whole harmonic diagonal is \textbf{anchor $+$ torsion forced}; $\delta=\tfrac12$ needs no
selection. \emph{Residue \mC:} $a_{13}=0$ and the gauge-orbit uniqueness are numerical (one component
across all seeds); promoting them to \mE\ via Gr\"obner on the RH constraints is the remaining
$U_{\mathrm{point}}$ step. \veri{v115\_anchor\_residue.py}


\begin{keybox}[The resonance theorem: $M_\infty=\mathbf1\iff a_{13}=0$ --- one gauge orbit \mE\ (\veri{v116\_resonance\_uniqueness.py})]
The anticipated Gr\"obner promotion turned out to be a \emph{linear} computation. Writing the
equivariant system at infinity as $dY/dw=-(1/w)(B_0+B_1w+\cdots)Y$ with
$B_m=\sum_k i^{km}\,U^kA_0U^{-k}$, the \emph{twisted} $\mu_4$ averages collapse \mE:
\[
  B_0=4\operatorname{diag}A_0,\qquad
  \boxed{\ B_1=4a_{12}E_{12}+4\overline{a_{13}}E_{31}\ }
\]
(only the eigenvalue-ratio $-i$ cells survive). The anchor exponents $\operatorname{diag}(2,1,1)$ have
resonance gap $1$; the level-$1$ formal-gauge equation $(1-\operatorname{ad}R_0)H_1=R_1$ is singular
exactly on the cells $\{(2,1),(3,1)\}$, where $B_1$ has entries $(0,\,4\overline{a_{13}})$; all higher
levels are non-resonant and the formal monodromy is $e^{-2\pi i\operatorname{diag}(2,1,1)}=\mathbf1$.
Hence \mE:
\[
  \boxed{\ M_\infty=\mathbf1\iff a_{13}=0\ }
\]
--- the transcendental Riemann--Hilbert condition collapses to \emph{one linear equation}.
\emph{Uniqueness corollary \mE:} $a_{13}=0$ $+$ the $(8,0,5)/144$ moduli $+$ phases $=$ diagonal gauge
$\Rightarrow$ the flat anchor locus is \textbf{exactly one gauge orbit}, the exact $A_0^\star$: within
the $\mu_4$-equivariant anchor class the $(U_{\mathrm{wall}})$ bundle datum is \emph{algebraically
unique}. \emph{Sharpness \mE:} $\|M_\infty-\mathbf1\|\sim8.5\,|a_{13}|$ ($0.17$ at $a_{13}=0.02$) vs
$10^{-10}$ at $A_0^\star$. \emph{Honest scope \mC:} the equivariant ansatz is the established $D_4$
selector input (\vref{v19}/\vref{v30}); the \emph{residue} side is now \mE, while the holonomy values (the harmonic
diagonal $(0,\tfrac12,\tfrac12)$) remain \mE\ --- transcendental monodromy of an exact system.
\end{keybox}

\begin{keybox}[The monodromy is the Weyl group of $A_3$ --- the holonomy is exact \mE\ (\veri{v117\_monodromy\_weyl\_a3.py})]
The last \mE\ item of the chain --- the holonomy \emph{values} --- closes: the unitarised monodromy of
the exact $A_0^\star$ system is itself \emph{exact}, with entries in $\tfrac12\Z[i]$:
\[
  \boxed{\ \widetilde M_0=\begin{pmatrix}0&-\tfrac{1+i}2&\tfrac{1-i}2\\[2pt]
  -\tfrac{1+i}2&-\tfrac i2&-\tfrac12\\[2pt]
  \tfrac{1-i}2&-\tfrac12&\tfrac i2\end{pmatrix}\ }
\]
unitary, $\det=1$, $\operatorname{tr}=0$, $\chi=\lambda^3-1$ (the cusp class \emph{exactly}),
$\widetilde M_0^3=\mathbf1$, and $\operatorname{diag}\widetilde M_0=(0,-\tfrac i2,+\tfrac i2)$ --- so
$d_1=0$ and $\delta=\tfrac12$ are now \emph{matrix entries}, not observations; $(\widetilde M_0U)^4=\mathbf1$
with $\operatorname{tr}(\widetilde M_0U)=1$ (the $\{1,i,-i\}$ branch, $d_1=0$ slice --- exactly where
\vref{v114}/\vref{v115} located the physical point). \emph{The group} \mE: exact enumeration of
$\langle U,\widetilde M_0\rangle$ gives order $24$ with order statistics $(1,9,8,6)$ and character
values $(3,-1,0,1)$ --- uniquely $S_4$ in its standard${}\otimes{}$sign representation, i.e.\ the
(twisted) reflection representation of the \textbf{Weyl group of the family lattice $A_3$}:
\[
  \boxed{\ \text{monodromy}(U_{\mathrm{wall}})=W(A_3)=S_4,\qquad |W(A_3)|=4!=24\ }
\]
--- the flavor wall realises the $A_3$ factor of the glue $D_5\oplus A_3+\mu_4$ as its Weyl group ($U$
$=$ a $4$-cycle, the $\mu_4$ deck; $\widetilde M_0$ $=$ a $3$-cycle, the family rotation).
\emph{Identification \mE:} the ODE monodromy of $A_0^\star$ equals the exact $\widetilde M_0$ to
$3.5\times10^{-10}$, and the invariant form $H$ is $U$-invariant to $10^{-11}$. \emph{Status:} holonomy
values \mE; only the analytic identification stays \mE. The $\delta$ thread
(\vref{v20} $\to$ \vref{v40}/\vref{v41} $\to$ \vref{v114} $\to$ \vref{v115} $\to$ \vref{v117}) is closed end to end.
\end{keybox}

\begin{keybox}[The hexagon is the sign-twisted family spectrum --- first exact H2 piece \mE\ (\veri{v118\_hexagon\_family\_dictionary.py})]
The missing ``$C_6\leftrightarrow\PP\setminus\mu_4$ dictionary'' (H2) acquires its first exact entry.
\emph{Sign-twist lemma \mE:} $-\widetilde M_0$ has order $6$ and
\[
  \boxed{\ \operatorname{spec}(\widetilde M_0)\,\cup\,\operatorname{spec}(-\widetilde M_0)=\mu_6\ }
\]
--- the six-site hypercharge hexagon spectrum \emph{is} the family monodromy spectrum together with its
sheet twist ($\Z_6=\Z_2\times\Z_3$ realised as $(-1)\times\langle\widetilde M_0\rangle$); the \vref{v20}
denominators $\tfrac54-\cos(r\pi/3)$ are exactly the eigen-denominators $|1-\zeta_6^r/2|^2$ of the two
family resolvents $(\mathbf1\mp\tfrac12\widetilde M_0)^{-1}$. \emph{Cyclotomic determinant \mE:}
$\det(\mathbf1-t\widetilde M_0)=1-t^3=(1-t)\Phi_3(t)$, so $\det(\mathbf1\mp\tfrac12\widetilde M_0)=
(\tfrac78,\tfrac98)$ --- and the lepton coefficients become resolvent determinants:
$c_e=|\mu_4|/(2\det_-)=\tfrac{16}7$, $c_\mu=\Nfam/(2\det_+)=\tfrac43$,
$c_\tau=|\mu_4|\det_-/\Nfam=\tfrac76$, product rule $=|\mu_4|/\det_+=\tfrac{32}9$. The lepton $7$ and
$9$ \emph{are} the family-resolvent determinants; $\delta=\tfrac12=|(\widetilde M_0)_{22}|$ is supplied
by the matrix itself. \emph{Sheet-extended group \mE\ (audit):} $\langle U,\widetilde M_0,-\mathbf1\rangle$
has order $48=\Omega_{\mathrm{adm}}=\Nfam\dim S^+$ --- the seed's quartic coefficient. \emph{Residue \mC:} the
dictionary's \emph{values} are exact; its \emph{address table} (which hexagon site each fermion
occupies; the quark composition) remains \vref{v20}'s established input / open --- nothing re-fished.
\end{keybox}
